%------------------------------------------------%
% Chap. 9 Using BNCmatmul from Python
%------------------------------------------------%
\chapter{\bi{Using BNCmatmul from Python}{PythonからのBNCmatmulの利用}}\label{chap:python}

\bi{This chapter shows how to call the DD/TD/QD (and DS/TS/QS) arithmetic,
FMA and elementary functions of BNCmatmul from Python via the standard
\texttt{ctypes} module, using the exported function API of
Sec.~\ref{sec:rdtqfunc} (\texttt{src/rdtq\_func.c} /
\texttt{include/rdtq\_func.h}) and the elementwise array functions of
Sec.~\ref{sec:bncelem}.  All code fragments in this chapter are taken from
\texttt{python/example\_rdtq.py}, which is shipped with the library and can
be executed as-is; \texttt{python/test\_rdtq.py} is the corresponding
regression driver.}{本章では、\ref{sec:rdtqfunc}節の関数エクスポートAPI
（\texttt{src/rdtq\_func.c} / \texttt{include/rdtq\_func.h}）と
\ref{sec:bncelem}節の要素毎の一括関数を使って、BNCmatmulの
DD/TD/QD（およびDS/TS/QS）演算・FMA・初等関数をPython標準の
\texttt{ctypes} モジュールから呼び出す方法を示します。本章のコード断片は
すべてライブラリに同梱の \texttt{python/example\_rdtq.py} から取ったもので、
そのまま実行できます。対応する回帰テストは \texttt{python/test\_rdtq.py}
です。}

\section{\bi{Building the shared library}{共有ライブラリのビルド}}

\bi{The exported symbols live in a small shared library
\texttt{python/librdtq.so}, built from \texttt{rdtq\_func.c} (99 scalar
functions) and \texttt{elem\_vector.c} (12 array functions) by the
legacy make targets:}{エクスポートされるシンボルは、\texttt{rdtq\_func.c}
（スカラー関数99個）と \texttt{elem\_vector.c}（一括関数12個）から
ビルドされる小さな共有ライブラリ \texttt{python/librdtq.so} に
収められます。レガシーmakeターゲットでビルドします:}

\begin{verbatim}
$ cd bncmatmul-0.24/src
$ make -f Makefile.legacy librdtq        # generic build
$ make -f Makefile.legacy librdtq_neon   # Arm NEON build
\end{verbatim}

\bi{or directly with the compiler:}{あるいはコンパイラを直接使う場合:}

\begin{verbatim}
$ gcc -O3 -ffp-contract=off -fPIC -shared -Iinclude \
      src/rdtq_func.c src/elem_vector.c -o python/librdtq.so -lm
\end{verbatim}

\bi{(\texttt{-ffp-contract=off} is mandatory, as everywhere in BNCmatmul.)
The full library shared objects \texttt{libbncmatmul-0.24*.so} contain the
same symbols plus the whole linear-computation API; \texttt{librdtq.so} is
merely the smallest self-contained subset for scripting use.}{（BNCmatmulの
他の部分と同様、\texttt{-ffp-contract=off} は必須です。）フルライブラリの
共有オブジェクト \texttt{libbncmatmul-0.24*.so} にも同じシンボルに加えて
線形計算API全体が含まれます。\texttt{librdtq.so} はスクリプト用途向けの
最小自己完結サブセットです。}

\section{\bi{Loading the library and printing full precision}{ライブラリのロードとフル精度表示}}

\bi{A DD/TD/QD value is an unevaluated sum of 2/3/4 doubles, so the natural
ctypes representation is a small \texttt{c\_double} array.  Since every limb
is an exact binary value, converting each limb to \texttt{decimal.Decimal}
and summing reproduces the full-precision value exactly --- Python's
\texttt{float} arithmetic would round it back to 53 bits:}{DD/TD/QD値は
2／3／4個のdoubleの未評価和なので、ctypesでは小さな \texttt{c\_double}
配列で表すのが自然です。各limbは厳密な2進値なので、limbごとに
\texttt{decimal.Decimal} へ変換して合計すればフル精度の値が厳密に
再現されます（Pythonの \texttt{float} 演算では53ビットに丸め戻されて
しまいます）:}

\begin{verbatim}
import ctypes as ct
from decimal import Decimal, getcontext

lib = ct.CDLL("./librdtq.so")

DD = ct.c_double * 2    # double-double  (~32 decimal digits)
TD = ct.c_double * 3    # triple-double  (~48 decimal digits)
QD = ct.c_double * 4    # quad-double    (~63 decimal digits)

getcontext().prec = 70
def to_decimal(limbs):
    return sum((Decimal(x) for x in limbs), Decimal(0))

x = DD(1.5, 0.0)        # x = 1.5 (high limb, low limb)
r = DD()
lib.rdd_exp(r, x)       # result first, as in rdd.h
print("dd exp(1.5) =", to_decimal(r))
# dd exp(1.5) = 4.48168907033806482260205546011928...
\end{verbatim}

\bi{Note the argument order: as in the \texttt{rdd.h} convention, the result
comes first.  The float-based families use \texttt{ct.c\_float * 2/3/4}
arrays and the \texttt{rds\_} / \texttt{rts\_} / \texttt{rqs\_}
prefixes.}{引数順に注意してください。\texttt{rdd.h} の慣習どおり、結果が
先頭です。float基のファミリは \texttt{ct.c\_float * 2/3/4} 配列と
\texttt{rds\_} / \texttt{rts\_} / \texttt{rqs\_} 接頭辞を使います。}

\section{\bi{Arithmetic, FMA and elementary functions}{算術演算・FMA・初等関数}}

\bi{The next fragment exercises quad-double square root, the certified
branch-free FMA (as an exact residual check $\sqrt{2}^2 - 2$), the
FMA-driven triple-double division, and the simultaneous
$\sin$/$\cos$:}{次の断片では、quad-double平方根、証明付きbranch-free FMA
（厳密な残差検査 $\sqrt{2}^2 - 2$ として）、FMA駆動triple-double除算、
$\sin$／$\cos$ の同時計算を使います:}

\begin{verbatim}
# quad-double: sqrt(2) to ~63 digits
two = QD(2.0, 0.0, 0.0, 0.0)
s = QD()
lib.rqd_sqrt(s, two)
print("qd sqrt(2) =", to_decimal(s))
# 1.41421356237309504880168872420969807856967187537694807317667973...

# residual sqrt(2)^2 - 2 via the certified FMA: fma(s, s, -2)
m2 = QD(-2.0, 0.0, 0.0, 0.0)
resid = QD()
lib.rqd_fma(resid, s, s, m2)
print("qd sqrt(2)^2 - 2 =", to_decimal(resid))   # ~0 (qd eps level)

# FMA-driven division: 1/3 in triple-double
a = TD(1.0, 0.0, 0.0);  b = TD(3.0, 0.0, 0.0);  q = TD()
lib.rtd_div_fma(q, a, b)
print("td 1/3 =", to_decimal(q))
# 0.33333333333333333333333333333333333333333333333327...

# simultaneous sin/cos (two results, then the input)
sin_x = DD();  cos_x = DD()
lib.rdd_sincos(sin_x, cos_x, x)

# div/sqrt-safe FMA: the scalar multiplier is a plain c_double
res = DD()
lib.rdd_fma_safe(res, DD(3.0, 0.0), ct.c_double(-0.5), DD(1.5, 0.0))
print("1.5 - 0.5*3 =", to_decimal(res))          # exactly 0
\end{verbatim}

\bi{Every exported function of Sec.~\ref{sec:rdtqfunc} is callable in the
same style: \texttt{rXX\_add} / \texttt{sub} / \texttt{mul} / \texttt{div} /
\texttt{sqrt}, \texttt{rXX\_fma} / \texttt{fma\_safe} / \texttt{div\_fma},
and \texttt{rXX\_exp} / \texttt{expm1} / \texttt{log} / \texttt{log10} /
\texttt{sin} / \texttt{cos} / \texttt{sincos} / \texttt{tan} /
\texttt{nint}.}{\ref{sec:rdtqfunc}節でエクスポートされる全関数が同じ流儀で
呼び出せます: \texttt{rXX\_add} / \texttt{sub} / \texttt{mul} /
\texttt{div} / \texttt{sqrt}、\texttt{rXX\_fma} / \texttt{fma\_safe} /
\texttt{div\_fma}、および \texttt{rXX\_exp} / \texttt{expm1} /
\texttt{log} / \texttt{log10} / \texttt{sin} / \texttt{cos} /
\texttt{sincos} / \texttt{tan} / \texttt{nint} です。}

\section{\bi{Elementwise array functions with NumPy}{NumPyと要素毎の一括関数}}

\bi{The elementwise array functions \texttt{bnc\_XX\_yy\_array}
(Sec.~\ref{sec:bncelem}) operate on limb-planar (SoA) arrays: limb $k$ of
all $n$ elements is one contiguous \texttt{double} array.  This maps
directly onto separate NumPy arrays, one per limb, passed as a
\texttt{double*[K]} pointer array --- no copying is involved, and one call
replaces $n$ round trips through ctypes:}{要素ごとの一括関数
\texttt{bnc\_XX\_yy\_array}
（\ref{sec:bncelem}節）はlimb-planar（SoA）配列を扱います。全$n$要素の
limb $k$ が1本の連続した \texttt{double} 配列です。これはlimbごとに
1本のNumPy配列をそのまま割り当て、\texttt{double*[K]} のポインタ配列として
渡す形に直接対応します（コピーは発生しません）。1回の呼び出しで、
ctypes経由の$n$回の往復を置き換えられます:}

\begin{verbatim}
import numpy as np

n = 100000
x_hi = np.linspace(-10.0, 10.0, n)     # high limbs
x_lo = np.zeros(n)                     # low limbs
r_hi = np.empty(n);  r_lo = np.empty(n)

def planar(*arrays):
    """double*[k] view of k numpy arrays (limb-planar layout)"""
    ptr = ct.POINTER(ct.c_double) * len(arrays)
    return ptr(*[a.ctypes.data_as(ct.POINTER(ct.c_double))
                 for a in arrays])

lib.bnc_dd_exp_array(ct.c_long(n),
                     planar(r_hi, r_lo), planar(x_hi, x_lo))

err = np.max(np.abs((r_hi + r_lo) / np.exp(x_hi) - 1.0))
print("max rel.err vs numpy =", err)   # ~2.2e-16 (binary64 comparison floor)
\end{verbatim}

\bi{The same pattern works for \texttt{bnc\_dd\_log\_array} /
\texttt{sin} / \texttt{cos} and, with three or four limb arrays, for the
\texttt{bnc\_td\_*\_array} / \texttt{bnc\_qd\_*\_array} functions.
The array layout is identical to the library's
\texttt{ddvector}/\texttt{tdvector}/\texttt{qdvector} types
(\texttt{vec->element}), so results can be exchanged with C code without
conversion.}{同じパターンが \texttt{bnc\_dd\_log\_array} /
\texttt{sin} / \texttt{cos} に、そしてlimb配列を3本または4本にすれば
\texttt{bnc\_td\_*\_array} / \texttt{bnc\_qd\_*\_array} にも使えます。
配列レイアウトはライブラリの
\texttt{ddvector}/\texttt{tdvector}/\texttt{qdvector} 型
（\texttt{vec->element}）と同一なので、C側との結果のやり取りに変換は
不要です。}

\section{\bi{Notes and caveats}{注意事項}}

\begin{itemize}
\item \bi{Result arguments come first (the \texttt{rdd.h} convention);
the C-level \texttt{c\_dd\_*} functions use the opposite (input-first)
order.  Mixing them up produces silently wrong results, so keep to the
\texttt{rXX\_*} API from Python.}{結果引数が先頭です（\texttt{rdd.h}の
慣習）。C側の \texttt{c\_dd\_*} 関数は逆（入力が先頭）です。取り違えると
静かに誤った結果になるため、Pythonからは \texttt{rXX\_*} APIを使って
ください。}
\item \bi{Scalar arguments must be wrapped explicitly
(\texttt{ct.c\_double(...)} / \texttt{ct.c\_float(...)}), e.g.\ the scalar
multiplier of \texttt{rXX\_fma\_safe}; bare Python floats passed by ctypes
default to \texttt{int}-like conversion and corrupt the call.}{スカラー引数は
明示的にラップしてください（\texttt{ct.c\_double(...)} /
\texttt{ct.c\_float(...)}）。例えば \texttt{rXX\_fma\_safe} のスカラー
被乗数です。素のPython floatはctypesの既定変換により呼び出しを壊します。}
\item \bi{Inputs are expected as normalized limb tuples.  Values
constructed as \texttt{(x, 0, ...)} from a Python float are always safe;
results returned by the library are normalized.}{入力は正規化された
limb組を想定しています。Python floatから \texttt{(x, 0, ...)} の形で
作った値は常に安全で、ライブラリが返す結果は正規化されています。}
\item \bi{For full-precision I/O beyond \texttt{to\_decimal}, the MPFR-based
conversion functions of \texttt{mpfr\_dtq\_sd.c} (e.g.\
\texttt{rdd\_get\_str}) are exported by the full
\texttt{libbncmatmul-0.24*.so} libraries and by \texttt{librdd.so}
(see \texttt{python/rdd.py}).}{\texttt{to\_decimal} 以上のフル精度入出力
には、\texttt{mpfr\_dtq\_sd.c} のMPFRベース変換関数（例:
\texttt{rdd\_get\_str}）がフルライブラリ \texttt{libbncmatmul-0.24*.so}
および \texttt{librdd.so} からエクスポートされています
（\texttt{python/rdd.py} を参照）。}
\end{itemize}
